% !TeX root = ../main.tex

\xchapter{面向长链条推理的平面几何基准构建}{A Plane Geometry Benchmark for Long-step Reasoning}

为了评估模型在复杂几何任务中的长链推理能力，本章构建了一个面向长链条平面几何推理的多模态评测基准。围绕这一目标，本章主要开展三方面工作：构建细粒度标注的长链条几何数据集，设计从原始图文到谓词逻辑表示的形式化工具链，并对多种代表性多模态大语言模型开展系统评测。

\xsection{引言}{Introduction}

在数学教育实践与几何自动求解研究中，平面几何问题的求解通常并非依靠单步公式套用即可完成，而是依赖多步、连续且可追溯的逻辑推导过程。每一步推理都需要严格依据几何定理与已知条件展开，并在前序结论的基础上逐步逼近目标结论。相较于短链推理任务，长链条推理对模型能力提出了更高要求：模型不仅需要具备几何知识调用能力，还需要在较长推理链条中持续保持条件约束，正确组织中间结论，并有效避免误差累积与推理路径偏移。这类能力要求也更加贴近中学阶段几何教学中“证明—推导—计算”交织的真实解题过程。

随着大语言模型和多模态大语言模型的发展，模型在基础几何问答与简单求解任务上已表现出一定能力，能够在有限推理深度下生成可用答案。然而，当题目涉及较长推理链条时，现有模型仍普遍存在推理断裂、多步传递过程中信息偏差等问题。造成这一现象的主要原因在于，现有几何数据基准对高难度、长链条推理场景的覆盖仍然不足，且多数评测仍以最终答案正确性为主，难以对中间推导过程进行细粒度刻画。因此，要系统评估模型在复杂几何任务中的真实能力，仅依赖短链样本或结果级评价是不充分的，亟需构建面向长链条平面几何推理的专用多模态数据基准。

基于上述背景，本章将面向长链条推理的数据基准构建任务定义为：从真实平面几何题目中筛选具有较长推理路径的问题样本，并围绕题目文本、几何图像、标准答案与推理过程建立统一的数据表示、标注规范和评测机制。该任务不仅要求数据样本具有较高的推理复杂度，还要求能够支持对图文条件理解、形式化表达、中间步骤合理性和最终结果正确性的综合评估。

围绕上述任务，本章构建了面向长链条平面几何推理的多模态评测基准 GeoLStep（Geometry Benchmark for Long-Step Reasoning）。本章的重点并不局限于传统意义上几何题目的收集与整理工作，而是在任务界定的基础上，系统设计适用于长链条推理场景的数据构建方案。针对平面几何问题图文信息交织、推理过程复杂的特点，本章建立了统一的图文联合标注与形式化表示框架，并进一步通过数据增强提升样本在推理深度与结构类型上的覆盖能力。与此同时，本章还从数据质量、统计特征与基线评测等多个层面对 GeoLStep 的基准有效性进行了验证，为后续长链条几何推理方法的研究与评估奠定了数据基础。

% -------------------------------------------------------------------
\xsection{GeoLStep 构建流程}{Construction of GeoLStep Benchmark}

\xsubsection{数据收集与处理}{Data Collection and Processing}

为构建面向真实考试场景的长链条平面几何推理基准 GeoLStep，本章系统收集了国内各省市 2020 年至 2024 年中考真题及代表性模拟试卷中的平面几何题目。相较于通用合成数据，此类题目在命题规范、知识点覆盖范围以及难度梯度设置等方面更具稳定性，因而能够为长链几何推理能力的评估提供更具代表性的测试分布。原始材料均来自互联网公开可获取的试卷资源。最终，共收集公开试卷 PDF 文件 383 套，并在此基础上对其中的平面几何题目进行系统筛选与整理。图~\ref{fig:paper_example} 给出了从原始试卷构建最终几何样本集合的流程示例。

该筛选流程主要包括两个核心环节。首先，在可求解性预筛阶段，为保证后续符号系统与形式化求解流程能够有效覆盖目标题目，本章基于题目结构对其可符号化建模能力进行预估，仅保留符合要求的欧氏几何问题，并剔除解析几何题目（如涉及坐标系与二次函数的问题）、以“可能的最大值/最小值”为求解目标的最值问题，以及其他符号系统难以直接建模的开放性问题等。

其次，在多题干处理阶段，针对包含多个设问的综合题，本章将题目主干文本与各小题文本按照逻辑顺序分别拼接，并拆分为若干独立样本。该处理保证了最终几何样本集合中的每个样本均具有完整自洽的题面和单一明确的求解目标。

\begin{figure}[H]
    \centering
    \includegraphics[width=0.9\textwidth]{figures/pdf_to_problem.pdf}
    \caption{试卷中选取平面几何题目样本示例}
    \label{fig:paper_example}
\end{figure}

基于上述数据来源与筛选方式，所构建数据集在题面信息组织方式与图像呈现风格上呈现出较为鲜明的特征。与部分公开几何数据集倾向于在图像中直接标注平行、垂直、角度数值等关键几何约束不同，国内中考试题的关键信息通常以自然语言形式完整给出，题目文本基本覆盖求解所需的主要条件；相应地，配图更多承担展示几何拓扑结构与作图对象的功能。因此，本数据集中的图像整体更为简洁规范，通常仅包含点、线段、圆和角弧等基础几何元素，整体以黑色线条为主，干扰信息较少，更贴近真实考试场景下的标准作图风格。在此基础上，中考几何题在设问形式和解题过程上通常也表现出较强的规范性，题目条件、推导路径与中间结论之间往往存在较为清晰的对应关系。这一特点有利于开展过程级标注、推理过程分析以及模型可解释性评测，也使本数据集更适合作为长链条平面几何推理研究的评测基础。

为了保证题面与解析的可用性，本章采用人工流程与自动化工具辅助相结合的方式完成数据整理与标注。整体标注框架如图~\ref{fig:c3_geolstep_annotation_framework} 所示。

\begin{figure}[htbp]
    \centering
    \includegraphics[width=0.95\textwidth]{figures/c3_geolstep_annotation_framework.pdf}
    \caption{GeoLStep 样本整理与统一标注框架}
    \label{fig:c3_geolstep_annotation_framework}
\end{figure}

具体而言，首先进行原始题目抽取，从 PDF 中对几何题目区域进行截图，分别保存高质量几何图像与对应的题面及解析区域；随后进行题目可求解性预筛，剔除题面缺失、图像不可辨识或解析信息不足的样本，以保证后续样本整理与标注工作的有效性。在文本识别与修正阶段，对题目文本与参考解答过程进行光学字符识别（Optical Character Recognition, OCR）获取文字信息，并使用多模态大语言模型（MLLM，如 GPT-4o\cite{hurst2024gpt4o}）进行内容校对与修正，在不改变题意的前提下消除 OCR 识别误差、补全漏识内容并统一文本表述。

在上述整理结果基础上，进一步完成求解过程、辅助线、难度信息标注和形式化转换。具体而言，基于修正后的参考解答，使用 MLLM 给出该题的推理步骤数以及是否需要作辅助线等标注信息，其中过程粒度遵循对“逻辑推理步骤”的定义；同时，在难度分级阶段，采用“易/中/难”三级标注，并依据是否需要辅助线帮助求解与推理长度进行划分，以反映解题的结构复杂度与推理深度。最后，对题目文本、几何图像与解答过程进行形式化解析，为后续图文形式化转换链路提供统一输入，再辅以人工复核进行校验，以保证最终结构化数据的质量与一致性。

\xsubsection{图文形式化转换链路}{Image-Text Formalization Pipeline}

为使几何题目能够被符号求解系统直接处理，并为过程核验、自动形式化证明等任务提供可验证的中间表示，有必要将自然语言描述转化为严格的形式化语言。近年来，相关研究普遍重视几何问题的形式化表达。例如，AlphaGeometry \cite{trinh2024alphageometry} 强调将自然语言条件转化为严密的逻辑构造，UniGeo \cite{chen2022unigeo} 也从统一多模态约束表示的角度对几何问题进行建模。这些工作表明，形式化语言不仅是连接自然语言与符号推理的关键桥梁，也是实现几何问题标准化表示的重要基础。

本章采用与 Inter-GPS \cite{lu2021intergps} 和 PGDPNet \cite{zhang2022pgdp} 等工作一致的谓词表示体系，如表\ref{tab:text_to_formal_example}中的一个典型几何题为例：将原始中文题目首先翻译为英文题面，随后通过规则化匹配或借助大语言模型的模板生成能力，将题干中“已知条件—几何关系—求解目标”逐条抽取并写成谓词表达（如 Equals、BisectsAngle、Parallel、Find 等）。这种表示方式能够把文本中相对松散、依赖上下文的叙述，显式拆解为可检查的对象与关系：例如，“直角”可等价为“角度为90°”，“平行”以线段对的形式给出；同理，长度相等、角相等、共线、点在圆上等常见条件约束也均可映射到统一的谓词空间。

\begin{table}[H]
    \centering
    \caption{题目文本到形式化语言的转换示例}
    \label{tab:text_to_formal_example}
    \small
    \begin{tabularx}{\textwidth}{@{} l X @{}}
        \toprule
        类型    & 内容                                                                                                                                                            \\
        \midrule

        题目文本  &
        如图，在$\triangle ABC$中，$AB=5$, $AC=9$，$AD$是$\angle BAC$的角平分线，点$E$是$BC$的中点，$EF//AD$，则$AF$的长是？                                                                            \\

        英文翻译  &
        In triangle $ABC$, $AB = 5$, $AC = 9$, $AD$ is the angle bisector of $\angle BAC$, $E$ is the midpoint of $BC$, $EF$ is parallel to $AD$, then the length of $AF$ is? \\

        形式化语言 &
        \begin{minipage}[t]{\linewidth}
            \small
            \setlength{\baselineskip}{0.9\baselineskip}
            Triangle(A,B,C)\\
            Equals(LengthOf(Line(A,B)), 5)\\
            Equals(LengthOf(Line(A,C)), 9)\\
            BisectsAngle(Line(A,D), Angle(B,A,C))\\
            IsMidpointOf(Point(E), Line(B,C))\\
            Parallel(Line(E,F), Line(A,D))\\
            Find(LengthOf(Line(A,F)))
        \end{minipage}                                                                                                                            \\

        \bottomrule
    \end{tabularx}
\end{table}

基于该谓词体系，形式化语言能够有效减少自然语言描述中的省略、歧义与表达变体问题，并将题目条件组织为结构化、可计算的约束集合，作为符号推理与过程核验的统一接口。在具体语法实现上，本章严格沿用 Inter-GPS 的符号化定义规范，以保证与主流几何求解器的兼容性，后续也可基于该体系进行不同形式化逻辑之间的转换。进一步地，本章构建了文本与图像两侧协同的自动形式化转换链路，分别面向题面语义的谓词生成与题图结构的形式化解析展开。

\subsubsection{题目文本形式化转换策略}

题目文本到形式化语言的映射本质上是一个高度结构化的语义解析任务。该任务要求模型严格遵循目标形式化语言的语法规范，对齐潜在的几何逻辑关系。为此，本研究以开源的 Qwen3-8B \cite{yang2025qwen3} 为基座，采用指令监督微调~\cite{long2022instructgpt}（Supervised Fine-Tuning）范式来训练专属的文本转换模型。相较于仅依赖通用提示词的大模型，面向任务的微调策略能够更稳定地学习几何对象、关系类型与目标谓词之间的映射规律，从而为后续批量形式化构建提供更可控的工具链支撑。

\paragraph{微调与测试数据集的构建}
高质量的指令数据是保证语言模型生成规范形式化语言的前提。本章的微调和测试数据集构建主要分为训练集扩充与独立测试集标注两个阶段。

在训练集构建方面，首先从开源数据集 Geometry3K \cite{lu2021intergps} 与 FormalGeo7k \cite{zhang2023formalgeo} 中提取英文题干及其对应的形式化输出，构造基础的“英文文本 $\rightarrow$ 形式化语言”监督样本对。其次，为增强模型对中文几何术语的理解能力及跨语言映射能力，将收集到的中文几何题干与上述英文样本进行混合对齐，以提升模型对中英文几何表达的统一建模能力。在此基础上，针对训练集引入了基于大语言模型的数据扩充策略，主要包括对原始中文题干进行同义重写、句式结构调整以及冗余条件注入等操作。其中，同义重写例如将“在 $\triangle ABC$ 中”等价改写为“已知三角形 $ABC$”；冗余条件注入仅作为表面表达层面的扰动，不引入会改变原有约束集合边界的伪几何条件，从而保证扩充样本与原始样本在形式化语义上的一致性。通过上述策略，有效微调样本规模扩充至原始数据的 5 倍，形成约 1.6 万条训练样本，以用于提升模型对复杂、非标准自然语言表述的泛化能力。

\begin{figure}[H]
    \centering
    \includegraphics[width=0.9\textwidth,keepaspectratio]{figures/fine_tune_data.pdf}
    \caption{微调数据集构建与数据扩充策略}
    \label{fig:fine_tune_data}
\end{figure}

测试集构建阶段从前文已构造的长链条数据集中独立采样，收集了约 1000 条几何题干作为测试数据，形成约 5500 条形式化语言样本。为保证评估标签的质量，测试集中的所有形式化语言均采用人工标注，并由具备相关学科背景的研究人员进行多轮交叉校验。此外，该测试集与微调所用的训练数据严格正交，排除了数据泄露的潜在风险。这批校验后的样本作为后续对比实验的基准数据（Ground Truth），为精确匹配率等核心量化指标的计算提供了事实依据。

\paragraph{微调流程}
微调流程采用监督微调范式。模型以自然语言题干作为输入，以目标形式化语言代码作为输出，通过最小化交叉熵损失（Cross-Entropy Loss）来优化模型权重，使其掌握跨语言几何关系转换与严格形式化语法生成能力。为兼顾显存效率与训练效果，实验引入 LoRA~\cite{hu2022lora}（Low-Rank Adaptation）参数高效微调技术，仅对低秩适配参数进行训练，而保持基座模型主体参数冻结。该策略能够在较低计算开销下完成领域适配，使模型学习几何题干到谓词表达之间的稳定映射关系。训练过程在单张 NVIDIA A800 (80GB) GPU 上完成。

\subsubsection{几何图像形式化解析流程}

在几何问题中，题目图像不仅起到辅助说明作用，还包含了大量题干文本中未显式给出的几何约束，例如点的共线关系、圆与直线的相交关系，以及由直角标记、平行标记等局部符号所表达的约束条件。因此，除了将题目文本转换为形式化表达之外，还需要进一步对几何图像进行形式化解析，从而获得完整的多模态结构化表示。

本实验采用 PGDPNet\cite{zhang2022pgdp} 作为几何图像解析工具，该模型面向平面几何图解析任务，能够对题图中的基础几何图元及其关系进行识别与建模，并输出结构化解析结果。利用该模型进行解析推理，可以将图像中的几何对象及其相互关系转换为图像侧的形式化语言。然而，PGDPNet 的解析过程高度依赖文本标签与局部符号信息的辅助，尤其需要较为准确的字符识别结果作为输入。考虑到通用 OCR 模型在几何题图场景下容易受到目标微小、图元紧邻以及符号遮挡等因素的干扰，本实验基于当前数据集构建了一个几何场景 OCR 训练集，并在此基础上对 PaddleOCR\cite{cui2025paddleocr} 进行微调，以针对性提升对象文本与几何符号的识别效果。整体流程如图 \ref{fig:image_parsing_pipeline} 所示。

\begin{figure}[htbp]
    \centering
    \includegraphics[width=0.9\textwidth]{figures/pipeline_flowchart.pdf}
    \caption{基于 OCR 微调与 PGDPNet 的几何图像形式化解析流程}
    \label{fig:image_parsing_pipeline}
\end{figure}

具体而言，首先使用半自动标注工具 PPOCRLabel~\footnote{https://github.com/PFCCLab/PPOCRLabel} 对本数据集中的几何图像进行细粒度标注，标注信息包括文本区域的边界框（BBox）坐标、对应文本内容，以及垂直、平行、角度等局部几何符号信息，如图 \ref{subfig:ocr_example} 所示。随后，在训练阶段以预训练的文本检测与识别模型作为初始参数，并利用上述标注数据对 OCR 模型进行微调，使其能够更好识别几何题图中的标识文本、垂直符号、平行符号以及角度标注等内容。微调结束后，将训练得到的 OCR 模型用于几何图像推理，获得题图中的文本边界框位置与识别结果，并将其输入 PGDPNet 框架。PGDPNet 进一步结合 OCR 结果与图像中的基础几何图元，对图像中的几何对象、符号标记及其空间关系进行解析，如图 \ref{subfig:pgdp_output} 所示，最终生成图像侧的形式化语言表示。

\begin{figure}[htbp]
    \centering
    \begin{subfigure}[b]{0.49\textwidth}
        \centering
        \includegraphics[width=\textwidth]{figures/c3_ocr_example.png}
        \caption{OCR 标注示例}
        \label{subfig:ocr_example}
    \end{subfigure}
    \hfill
    \begin{subfigure}[b]{0.49\textwidth}
        \centering
        \includegraphics[width=\textwidth]{figures/c3_pgdp_output.png}
        \caption{PGDPNet 解析结果示例}
        \label{subfig:pgdp_output}
    \end{subfigure}

    \caption{几何图像 OCR 标注与 PGDPNet 解析示例}
    \label{fig:ocr_pgdp_example}
\end{figure}

在获得自动解析结果后，实验中进一步对其进行人工校验与修正。校验内容主要包括点名称、线段名称、圆名称等几何对象识别结果，以及由这些对象和关系生成的图像侧形式化语言。通过该过程修正 OCR 误识别、几何对象漏检或关系解析错误等问题，最终为整个数据集补充较为可靠的几何图像形式化标注。

% -------------------------------------------------------------------
\xsection{数据统计与质量分析}{Data Statistics and Quality Analysis}

\xsubsection{数据统计特征}{Statistical Characteristics}

经过严格筛选与清洗后，本章最终构建了一个长链条几何题目为主的数据基准，
% ，其中计算题 1274 道，证明题 823 道。
数据集中每条样本均由题目文本和对应几何图像共同构成，属于典型的图像--文本多模态数据。图~\ref{fig:data_anna_example} 展示了一个几何题目样本及其结构化标注示例。标注信息主要涵盖三个部分：首先是题目元数据，包括题目类型、数值结果、难度等级、辅助线类型、候选选项及单位信息等；其次是语义逻辑形式，包括图形逻辑表达与文本逻辑表达；最后是求解过程参考信息，包括参考求解过程、步骤数以及分步解释等。

\begin{figure}[H]
    \centering
    \includegraphics[width=0.8\textwidth]{figures/data_anna_example.pdf}
    \caption{几何题目样本的结构化标注示例}
    \label{fig:data_anna_example}
\end{figure}

题目类型、辅助线需求和难度等级三个维度的分布情况如图~\ref{fig:c3_dataset_distribution} 所示，图中能直观呈现 GeoLStep 的任务构成与样本复杂性。可以看到，数据集同时包含计算题与证明题，并覆盖一定比例在参考解法中需要构造辅助线的样本；在难度方面，简单题、中等题和难题均有分布。整体来看，GeoLStep 的样本结构较为丰富，可用于考察模型在几何理解、构造推理和复杂问题求解等方面的能力。

\begin{figure}[htbp]
    \centering

    \begin{subfigure}[b]{0.32\textwidth}
        \centering
        \includegraphics[width=\textwidth]{figures/c3_problem_type_distribution.pdf}
        \caption{题目类型分布}
        \label{subfig:c3_problem_type_distribution}
    \end{subfigure}
    \hfill
    \begin{subfigure}[b]{0.32\textwidth}
        \centering
        \includegraphics[width=\textwidth]{figures/c3_auxiliary_line_distribution.pdf}
        \caption{辅助线需求分布}
        \label{subfig:c3_auxiliary_line_distribution}
    \end{subfigure}
    \hfill
    \begin{subfigure}[b]{0.32\textwidth}
        \centering
        \includegraphics[width=\textwidth]{figures/c3_difficulty_distribution.pdf}
        \caption{难度等级分布}
        \label{subfig:c3_difficulty_distribution}
    \end{subfigure}

    \caption{GeoLStep 数据集样本类型分布}
    \label{fig:c3_dataset_distribution}
\end{figure}

本章采用参考解中的逻辑推理步骤数来统计推理链长度。求解步数分布如图~\ref{fig:dataset_steps_count} 所示，整体主要集中在 $3\sim10$ 步区间，平均推理步数约为 7.44 步。数据集中，求解过程包含 $1\sim4$ 个推理步骤的题目共 509 道，占比约 24.3\%；包含 5 步及以上推理步骤的题目共 1588 道，占比约 75.7\%，构成本基准重点关注的长链推理样本。其中，$5\sim12$ 步题目占比最高，为 63.6\%；$13\sim26$ 步的高复杂度长链样本占比约 12.1\%。该分布表明，GeoLStep 在保留基础短链题目的同时，以中长链题目为主体，并覆盖一定数量的复杂连续推理样本，可用于考察模型在不同推理深度下的几何推断与定理使用能力。

保留一定比例的短步骤题目主要基于以下两方面考量。首先，保持难度分布的完整性与数据的真实性。在真实的几何问题分布中，简单题也是重要组成部分。其次，作为基础能力的对照组。如果模型在短步骤题目上表现不佳，说明其基础几何概念理解存在缺陷；若基础题表现良好但在长链条题上显著下降，则能更精准地定位其多步推理能力的瓶颈。最后，这也支持了渐进式推理的研究，有助于探索模型从浅层推理向深层推理迁移的行为模式。

\begin{figure}[H]
    \centering
    \includegraphics[width=0.9\textwidth]{figures/dataset_steps_count.png}
    \caption{数据集题目推理步骤数分布}
    \label{fig:dataset_steps_count}
\end{figure}

基于前述数据集统计，下文介绍图文形式化表示的构建结果。考虑到当前符号系统尚不能有效支持证明题求解过程的完整形式化与自动验证，所以本章将最终答案为数值的计算题作为文本形式化转换的主要对象。基于前文所述的微调方法与 Qwen3-8B~\cite{yang2025qwen3}，并结合人工标注和校验，本章构建了计算题文本侧的形式化表示。为保证表达结果的规范性与可执行性，所有文本形式化语言均通过 Inter-GPS~\cite{lu2021intergps}、GeoDRL~\cite{peng2023geodrl} 等几何求解方法的文本解析器测试，确保其能够被正确解析。

表~\ref{tab:formalization_statistics} 汇总了图文两侧形式化标注的整体统计结果。文本侧主要覆盖题目中的已知条件、求解目标及相关数量约束；图像侧则主要用于补充几何图形中的结构信息。由于中考几何题的核心数量条件通常已在题目文本中明确给出，图像侧提取结果更多集中于点、线、圆以及垂直、平行等基础拓扑关系。经统计，数据集中图像总共解析得到点 10988 个、线 20194 条、圆 747 个。

\begin{table}[htbp]
    \centering
    \caption{GeoLStep 图文形式化标注统计}
    \label{tab:formalization_statistics}
    \small
    \begin{tabular*}{\textwidth}{@{\extracolsep{\fill}}lcc}
        \toprule
        统计项 & 文本侧形式化 & 图像侧形式化 \\
        \midrule
        样本数量 & 1274 道计算题 & 1673 张图像 \\
        形式化语句总数 & 7389 条 & 11006 条 \\
        平均每题/每图语句数 & 5.80 条 & 6.58 条 \\
        主要信息类型 & 主要条件、求解目标 & 点线圆拓扑、垂直平行关系 \\
        \bottomrule
    \end{tabular*}
\end{table}

针对文本侧形式化语言，其表达式分布与谓词类型统计结果如图~\ref{fig:text_logic_analysis} 所示。可以看到，单题形式化表达式条数主要集中在中等规模区间，如 $3\sim8$ 条，表明多数计算题在形式化表示层面具有相对稳定的条件规模，通常能够通过有限数量的表达式刻画题目中的已知条件与求解目标。从高频谓词分布来看，长度关系和求解目标相关谓词占比较高，同时垂直、角度和线段相等等关系也较为常见。这一现象表明，GeoLStep 中的计算题主要围绕线段长度、角度关系和基本位置关系展开，文本侧形式化表示能够较好覆盖初中几何计算题中常见的数量关系与结构约束。

\begin{figure}[htbp]
    \centering
    \begin{subfigure}[b]{0.49\textwidth}
        \centering
        \includegraphics[width=\textwidth]{figures/figure_a_expression_counts.pdf}
        \caption{形式化表达式条数分布}
        \label{subfig:expr_count_dist}
    \end{subfigure}
    \hfill
    \begin{subfigure}[b]{0.49\textwidth}
        \centering
        \includegraphics[width=\textwidth]{figures/figure_b_predicate_frequencies.pdf}
        \caption{最高频谓词分布}
        \label{subfig:predicate_dist}
    \end{subfigure}

    \caption{几何计算题文本形式化语言统计}
    \label{fig:text_logic_analysis}
\end{figure}

表~\ref{tab:image_predicates} 给出了图像侧主要谓词类型及其数量分布统计。可以看到，点线关系和点圆关系占据主要部分，说明图像解析结果主要用于恢复几何图形的基本拓扑结构；相比之下，垂直、平行和角度数值等关系通常已在题面文本中明确描述，图像侧仅对其中可由图形标记识别的部分进行补充提取，因此其数量相对较少。整体而言，图像信息主要承担结构约束补充与拓扑关系呈现功能，核心数量条件仍主要来自题目文本，这与中考试题图文协同的命题特点一致。

\begin{table}[htbp]
    \centering
    \caption{图像侧形式化语言主要谓词分布统计}
    \label{tab:image_predicates}
    \small
    \begin{tabular*}{\textwidth}{@{\extracolsep{\fill}}clc@{}}
        \toprule
        谓词含义 & 形式化语言表示 & 数量（条） \\
        \midrule
        点在线上 & \texttt{PointLiesOnLine(Point, Line)} & 7220 \\
        点在圆上 & \texttt{PointLiesOnCircle(Point, Circle)} & 2843 \\
        垂直关系 & \texttt{Perpendicular(Line, Line)} & 568 \\
        平行关系 & \texttt{Parallel(Line, Line)} & 216 \\
        角度数值 & \texttt{Equals(MeasureOf(Angle), Number)} & 159 \\
        \bottomrule
    \end{tabular*}
\end{table}

\xsubsection{形式化转换性能分析}{Performance Analysis of Formalization Conversion}

为验证前述图文形式化转换链路的有效性，本节分别从文本到形式化语言转换和图像 OCR 解析两个方面，对工具链的性能进行对比分析。前者关注模型是否能够稳定生成符合语法且语义完整的谓词表达，后者关注 OCR 输入质量对几何图像形式化解析结果的影响。

\subsubsection{文本转换实验结果}

为客观评估模型将自然语言转换为形式化语言的质量，实验使用语法合法率（Syntax Valid Rate, SVR）、约束召回率（Constraint Recall, CR）与精确匹配率（Exact Match, EM）三项核心指标。其具体定义如下：

语法合法率（SVR）：衡量模型生成的形式化语言表达式是否符合目标的语法规范，能够被文本解析器\cite{lu2021intergps}成功编译而无语法错误。
\begin{equation}
    \mathrm{SVR} = \frac{N_{\mathrm{valid}}}{N_{\mathrm{total}}} \times 100\%,
\end{equation}
其中，$N_{\mathrm{total}}$ 为测试集样本总数，$N_{\mathrm{valid}}$ 为语法合法的样本数。

约束召回率（CR）：评估模型对原始题目中关键几何条件与求解目标的覆盖程度。具体而言，本文将每条形式化语句视为一个独立约束项，并在谓词类型及其参数对象均匹配时认为该约束被成功召回。CR 统计所有样本中被正确生成的真实约束项占真实约束项总数的比例，因此能够反映模型是否遗漏题目中的重要条件，但不要求整道题目的形式化结果与标注完全一致。
\begin{equation}
    \mathrm{CR} = \frac{\sum_{i=1}^{N_{\mathrm{total}}} |C_{\mathrm{pred}}^{(i)} \cap C_{\mathrm{gt}}^{(i)}|}{\sum_{i=1}^{N_{\mathrm{total}}} |C_{\mathrm{gt}}^{(i)}|} \times 100\%,
\end{equation}
其中，$C_{\mathrm{pred}}^{(i)}$ 和 $C_{\mathrm{gt}}^{(i)}$ 分别表示第 $i$ 个样本中模型预测的约束集合与人工标注的真实约束集合。

精确匹配率（EM）：要求模型生成的完整形式化结果与真实标签完全一致，或在抽象语法树上逻辑等价，即只有当一个样本中的所有约束项均被正确生成且不存在额外错误约束时，才记为匹配成功，是衡量下游符号求解可靠性的最严格标准。
\begin{equation}
    \mathrm{EM} = \frac{N_{\mathrm{match}}}{N_{\mathrm{total}}} \times 100\%,
\end{equation}
其中，$N_{\mathrm{match}}$ 为实现精确匹配的样本数量。

为全面评估文本形式化转换模型的有效性，本章设置了两类对比基线：第一类是基于少样本提示（Few-shot Prompt）的通用大模型，这类模型不进行权重更新，仅依赖上下文学习能力来捕捉形式化语言的映射规律；第二类是同等参数规模（约 8B--9B）且经过相同特定任务微调的开源模型，用于验证不同基座模型在吸收垂直领域知识后的表现差异。

各模型在测试集上的具体评估结果如表 \ref{tab:text_results} 所示。从整体对比结果可以看出，基于提示词的开源小模型（Llama-3.1-8B、GLM-4-9B、Qwen3-8B）在三项指标上整体偏低，其精确匹配率仅分布在 60.5\% 至 64.2\% 之间，尚不足以稳定支撑下游的符号求解或语义理解。相比之下，在少量样本提示的辅助下，闭源模型 GPT-4o mini、GPT-4o 以及大参数开源模型 DeepSeek-V3 均表现出更强的结构化生成能力，说明大规模模型在上下文模式归纳以及语法约束遵循方面具有明显优势。

\begin{table}[htbp]
    \centering
    \caption{不同模型在文本到形式化语言转换任务上的性能对比（\%）}
    \label{tab:text_results}
    \small
    \renewcommand{\arraystretch}{0.85}
    \setlength{\tabcolsep}{6pt}
    \begin{tabular*}{\textwidth}{@{\extracolsep{\fill}}lccc}
        \toprule
        模型 & 语法合法率 & 约束召回率 & 精确匹配率 \\
        \midrule

        \multicolumn{4}{@{}l}{\small\textit{Few-shot Prompt-based}} \\
        Llama-3.1-8B\cite{grattafiori2024llama}                & 80.1 & 73.2 & 60.5 \\
        GLM-4-9B\cite{glm2024chatglm}                          & 81.3 & 74.7 & 62.4 \\
        Qwen3-8B\cite{yang2025qwen3}                           & 82.6 & 75.6 & 64.2 \\
        GPT-4o mini\cite{hurst2024gpt4o}                       & 92.1 & 89.7 & 84.5 \\
        DeepSeek-V3\cite{liu2024deepseek}                      & 92.4 & \underline{90.5} & 85.9 \\
        \textbf{GPT-4o\cite{hurst2024gpt4o}}                   & \textbf{94.5} & \textbf{92.8} & \textbf{90.2} \\

        \cmidrule{1-4}

        \multicolumn{4}{@{}l}{\small\textit{Fine-tuning-based (\textasciitilde 8B--9B)}} \\
        Llama-3.1-8B$^*$\cite{grattafiori2024llama}            & 89.2 & 86.1 & 82.3 \\
        GLM-4-9B$^*$\cite{glm2024chatglm}$^*$                  & 91.5 & 87.8 & 83.6 \\
        \textbf{Qwen3-8B$^*$\cite{yang2025qwen3} (Ours)$^*$}   & \underline{92.8} & 90.2 & \underline{86.1} \\

        \bottomrule
    \end{tabular*}

    \vspace{1ex}
    \raggedright
    \footnotesize 注：$^*$表示经过本文特定数据集微调后的模型。
    \normalsize
\end{table}

进一步观察微调组可知，经过特定数据与 LoRA 微调后，开源小模型在文本到形式化语言转换任务上的性能均显著提升。其中，Qwen3-8B 的语法合法率、约束召回率与精确匹配率较未微调前分别提升了 10.2\%、14.6\% 与 21.9\%，并在语法合法率（92.8\%）与精确匹配率（86.1\%）上超过 GPT-4o mini 和 DeepSeek-V3，约束召回率（90.2\%）也与 DeepSeek-V3 的 90.5\% 基本接近。虽然其整体性能仍低于最优的 GPT-4o，但闭源模型受限于可控性与本地部署条件，DeepSeek-V3 等大参数模型部署时约需 680GB 文件空间；相比之下，微调后的 Qwen3-8B 仅需约 20GB 容量，更便于本地部署和系统集成，也更适合面向几何形式化转换这一专门任务进行优化。

\subsubsection{图像 OCR 解析性能对比}

为了量化评估 OCR 质量对下游图像形式化解析的影响，本实验在统一的 PGDPNet 解析框架下，对比了不同 OCR 输入来源下的解析性能。具体而言，对比基准包含四类：人工标注、微调后的 PaddleOCR（PaddleOCR-FT）、原生 PaddleOCR 以及通用识别大模型 DeepSeekOCR~\cite{wei2025deepseekocr}。以经人工严格校验后的图元（点、线、圆）及形式化语言结果为真实基准，统计各方法的正确率；同时，计算各方法预测结果中误检与冗余实例占其输出实例总数的比例，以此评估多余率。对比结果如图~\ref{fig:ocr_performance} 所示。

\begin{figure}[H]
    \centering
    \begin{subfigure}[b]{0.49\textwidth}
        \centering
        \includegraphics[width=\textwidth]{figures/ocr_accuracy.pdf}
        \caption{各方法识别正确率对比}
        \label{subfig:ocr_accuracy}
    \end{subfigure}
    \hfill
    \begin{subfigure}[b]{0.49\textwidth}
        \centering
        \includegraphics[width=\textwidth]{figures/ocr_excess_rate.pdf}
        \caption{各方法识别多余率对比}
        \label{subfig:ocr_excess_rate}
    \end{subfigure}

    \caption{不同 OCR 模型在各几何图元及形式化语言上的识别性能对比}
    \label{fig:ocr_performance}
\end{figure}

可以观察到，通用 OCR 模型在几何图表场景下存在明显瓶颈：原生 PaddleOCR 与 DeepSeekOCR 在形式化语言层面的正确率仅为 32\% 和 11\%，且多余率分别达到 44\% 和 45\%。这表明，通用 OCR 容易受到几何题图中微小目标、紧邻图元及复杂符号的干扰，从而影响后续 PGDPNet 的关系解析。相比之下，经过几何场景数据微调的 PaddleOCR-FT 性能显著提升，其在点和形式化语言上的正确率分别达到 83\% 和 76\%，并将形式化语言的多余率降至 23\%。尽管受限于几何图像本身较高的解析难度，自动方法在复杂图元上的表现仍与人工提供 OCR 输入时存在一定差距，但该对比实验说明了领域垂直训练的有效性。微调后的模型能够更好适应几何图像的分布特征，在保证基本解析精度的同时，显著降低后续构建同类大规模数据集时对人工 OCR 标注与校验的依赖。

\xsubsection{数据质量分析}{Data Quality Analysis}

GeoLStep 不仅包含题目文本与几何图像，还进一步提供了参考解答、推理步骤数、辅助线信息以及图文两侧的形式化表示。因此，本章的数据质量不仅体现在原始题目的正确性上，还体现在多模态信息、过程级标注与形式化结果之间的一致性上。为保证该基准能够稳定支撑后续的过程评估与模型对比，本节从原始题目质量、结构化标注一致性、图文形式化结果可靠性以及数据划分独立性四个方面进行分析。

\paragraph{原始题目质量控制}
本章数据来源于真实中考试卷与代表性模拟试卷。相较于网络问答平台或全自动合成题目，此类题目通常已经过命题、排版、校对和答案审阅等环节，题面规范性较强，且知识点覆盖和难度梯度较为稳定。在原始样本筛选阶段，本章进一步剔除了题面缺失、图像模糊、解析残缺、目标不明确或不适宜进行符号建模的样本，从源头上保证了题目内容的完整性与可求解性。对于综合题，则通过主干与子问拼接的方式将其拆分为独立样本，避免同一样本中存在多个目标混杂的问题，从而提升后续结构化标注与评测的一致性。

\paragraph{结构化标注一致性控制}
在样本整理完成后，本章对题目类型、最终答案、难度等级、辅助线需求、参考解答过程和推理步骤数等信息进行了统一标注。由于长链条几何题的关键难点在于连续推导过程，因此本章将“逻辑推理步骤”而非公式行数或文字分句作为步骤统计粒度，以更贴近真实求解中的关键中间结论传递。对于辅助线信息，则依据参考解答中是否显式构造新几何对象进行判断。整体上，标注过程采用“自动辅助生成—人工修订—统一复核”的方式完成：先由 OCR 与 MLLM 工具生成初始文本与步骤信息，再由人工结合原题、配图和参考解析进行修订，最后对存在歧义或标注不一致的样本进行统一复核。该流程在保证标注效率的同时，尽可能降低了单次自动识别或主观判断带来的偏差。

\paragraph{图文形式化结果可靠性控制}
考虑到图文形式化结果直接服务于后续的过程评估与符号推理，本章对文本侧与图像侧的形式化转换均设置了专门的质量控制环节。对于文本侧形式化结果，一方面通过语法合法性检查过滤括号不匹配、谓词名称错误、参数缺失等格式问题，另一方面通过人工校验确认谓词集合是否准确覆盖题干中的关键已知条件与待求目标。对于图像侧形式化结果，则先利用 OCR 与 PGDPNet 自动生成候选图元关系，再针对点名识别、图元漏检、关系误判等常见错误进行人工修正。前一小节中的对比实验也进一步说明，经过面向几何场景的微调后，文本转换模型与 OCR 模型均能够取得较高的结构化转换质量，从而为大规模构建过程中的自动化处理提供可靠支撑。

\paragraph{数据划分与泄漏控制}
为保证后续评测结果的可信性，本章在工具链训练、测试以及最终数据集整理过程中，注意保持不同用途样本之间的独立性。文本形式化转换任务所使用的测试集与微调训练数据严格正交，避免模型通过记忆训练样本而导致评估结果虚高。与此同时，最终 GeoLStep 数据集中保留了原始题面、配图、参考解析与结构化标注之间的对应关系，便于后续对异常样本进行追溯与再次核验。这样的组织方式既有利于保证当前基准构建的严谨性，也为后续数据扩展和误差分析预留了空间。

综合来看，GeoLStep 的数据质量控制覆盖了从原始题目筛选、结构化标注、图文形式化转换到评测用数据独立划分的完整流程。该流程保证了数据集既保留真实教育场景中的题目复杂性，又具备支持过程级评估和符号化分析所需的结构化一致性，为后续基线实验和长链条几何推理方法研究提供了较为可靠的数据基础。

% -------------------------------------------------------------------
\xsection{评估协议}{Evaluation Protocol}

\xsubsection{评估指标}{Evaluation Metrics}
% 数据集按 6:2:2 的比例随机划分为训练集、验证集和测试集。为了在防止数据泄露的同时保证各子集难度分布一致，本章采用了基于求解步长的分层采样，并进行了严格的交叉去重。

传统几何题评测通常依赖人工评审或者规则匹配最终结果，但仅依据最终答案正确率又难以反映模型在中间推理过程中可能出现的逻辑错误。因此在指标设计上，本章定义了答案准确率（$\mathrm{Acc}_{\text{ans}}$）和过程逻辑正确率（$\mathrm{Acc}_{\text{step}}$）来评估模型求解能力：

答案准确率考察最终输出结果匹配程度。计算题的正确性基于数值容差判定，证明题则取决于结论是否与参考答案一致：
\begin{equation}
    \mathrm{Acc}_{\text{ans}}=\frac{N_{\text{correct}}}{N}\times 100\%
    \label{eq:ans_acc},
\end{equation}
其中，$N$ 表示题目总数，$N_{\text{correct}}$ 表示答案预测正确的题目数。

过程逻辑正确率用于量化模型推理逻辑的严密性。一个生成的步骤必须同时满足前提充分、定理引用恰当且演绎正确，才能被视为有效：
\begin{equation}
    \mathrm{Acc}_{\text{step}}=\frac{N_{\text{valid}}}{N_{\text{all}}}\times 100\%,
    \label{eq:step_acc}
\end{equation}
其中，$N_{\text{all}}$ 为总步骤数，$N_{\text{valid}}$ 为有效步骤数。

\xsubsection{评估机制}{Evaluation Mechanism}

在结果评判阶段，为确保评估的严谨性，答案准确率（$\mathrm{Acc}_{\text{ans}}$）的计算采用正则提取、模型校验的双重机制：首先提取模型输出的最终答案，并与标准答案进行硬匹配；为避免因数学表达格式差异导致的误判，进一步引入评估模型对匹配失败的样例进行二次等价性复核。在评测设置中，本章以参考解的推理步骤数对题目进行难度分层，用于近似刻画推理链长度，并据此分析模型在不同长链条区间下的性能变化。

过程逻辑正确率计算上，基于前文数据集提供的参考求解过程与解释标注信息，本章使用自动化的过程评估机制，引入多模态大模型（如OpenAI o4-mini~\cite{openai2025o3o4mini}、Gemini-2.5-Pro~\cite{comanici2025gemini25}）作为主要评估器。基本思路为：通过设计结构化评测提示模板，对待测模型生成的每一个推理过程进行逐步审查，并要求评估器在欧氏几何公理体系与题目已知条件约束下作出判断，指出当前步骤是否正确。

具体而言，对于任意步骤 $s_i$，评估过程主要考察前提合法性、定理适用性与推导正确性三个方面，即判断该步骤所依赖的几何条件是否已经由题干文本、图像信息或前序推理过程给出，判断该步骤所引用的几何定理、公理或性质是否适用于当前拓扑结构与推理上下文，以及判断由该步骤得到的中间结论在逻辑关系与数值计算上是否正确。

为进一步比较不同评估模型在过程评估任务中的可靠性，本章随机抽取了 400 个推理样本，并以 OpenAI o3~\cite{openai2025o3o4mini} 的求解结果作为统一评估对象，组织人工专家对其推理过程与最终答案进行了校验，形成参考标注。在此基础上，分别使用 Claude-4.5~\cite{anthropic2025sonnet45}、Gemini-2.5-Pro~\cite{comanici2025gemini25}、GPT-4o~\cite{hurst2024gpt4o} 和 OpenAI o4-mini~\cite{openai2025o3o4mini} 对同一批样本进行自动评估，并将其判断结果与人工专家标注进行对比，统计步骤层面与答案层面的一致性，结果如表 \ref{tab:evaluator_comparison} 所示。可以看出，o4-mini 在步骤评估一致率和答案评估一致率上分别达到 97.6\% 和 99.58\%，均优于其他候选评估模型，说明其自动评估结果与人工判断最为接近。因此，后续实验中的过程逻辑正确率评估统一采用 o4-mini 作为自动评估器。

\begin{table*}[htbp]
    \centering
    \small
    \caption{不同评估模型与人工专家标注结果的一致性对比（\%）}
    \label{tab:evaluator_comparison}
    \begin{tabularx}{\textwidth}{XYY}
        \toprule
        评估模型                                       & 步骤评估一致率        & 答案评估一致率        \\
        \midrule
        Claude-4.5~\cite{anthropic2025sonnet45}    & 91.88          & 96.67          \\
        GPT-4o~\cite{hurst2024gpt4o}               & 91.95          & 97.08          \\
        Gemini-2.5-Pro~\cite{comanici2025gemini25} & 93.70          & 98.75          \\
        \textbf{o4-mini~\cite{openai2025o3o4mini}} & \textbf{97.60} & \textbf{99.58} \\
        \bottomrule
    \end{tabularx}
\end{table*}


% -------------------------------------------------------------------
\xsection{基线实验}{Baseline Experiments}

\xsubsection{实验设置}{Experimental Setup}

\subsubsection{评测设置}
为保证评估结果的公平性与可比性，本章对不同模型的部署方式、提示策略及解码参数等设置进行了统一控制。

在部署方式方面，几何专用模型均基于其官方代码仓库与推荐配置在本地复现，并在两张 NVIDIA A800（80GB）GPU 上完成部署和推理。对于闭源通用大模型，实验中统一采用远程 API 方式调用；对于开源通用模型，则按照官方推荐环境进行本地部署与推理。

在模型输入方面，部分几何专用模型受到输入形式的限制，无法直接处理原始图像与自然语言题干，因此实验中使用的是前文构建的计算题测试子集及其配套的形式化表示资源，并直接将其中预先标注的形式化语言作为模型输入，而不额外引入图文解析预处理模块。这样可以尽量避免视觉识别与文本解析误差对后续求解过程的干扰，使评估结果更集中地反映模型本身的几何推理与求解能力。基于这一设定，本章对这部分几何专用模型仅统计最终求解准确率，而不进一步评估其中间推理过程的正确性。相比之下，通用大语言模型均直接接收原始几何图像与题干文本作为输入，以尽可能保留题目的原始信息，并减少中间表示转换带来的额外偏差。

在提示策略上，所有通用大语言模型均采用零样本（Zero-shot）设置，不提供示例，仅通过任务指令约束输出格式与推理过程。采用这一设置主要有两方面考虑：一方面，少样本（Few-shot）提示容易引入额外的上下文增益，使模型性能受到示例选择与示例质量的影响；另一方面，零样本评测更符合本工作关注模型原生长链条推理能力的目标。提示词（Prompt）要求模型依据题目给定信息逐步推导，明确指出所使用的几何性质、定理或等量关系，并按照统一格式输出最终答案。本章采用的英文与中文提示词模板如图~\ref{fig:prompt_example}所示。

\begin{figure}[H]
    \centering
    \includegraphics[height=5.8cm]{figures/prompt_example.pdf}
    \caption{大语言模型提示词示例}
    \label{fig:prompt_example}
\end{figure}

为进一步降低不同 API 接口默认配置带来的随机性干扰，本章统一设定了大模型调用的核心生成控制参数，对于部分接口参数定义不完全一致的模型，遵循“语义对齐、数值相等”的原则进行调整。具体配置如表~\ref{tab:reasoning_params} 所示。

\begin{table}[H]
    \centering
    \small
    \caption{核心生成控制参数设置表}
    \label{tab:reasoning_params}
    \renewcommand{\arraystretch}{0.85}
    \begin{tabularx}{\textwidth}{YYY}
        \toprule
        参数名称               & 参数值    & 参数说明    \\
        \midrule
        n                  & $1$    & 候选返回数量  \\
        seed               & $42$   & 随机种子    \\
        top\_p             & $0.95$ & 核采样阈值   \\
        temperature        & $0.2$  & 生成随机性控制 \\
        max\_tokens        & $2048$ & 最大输出词元数 \\
        frequency\_penalty & $0.0$  & 频率惩罚系数  \\
        \bottomrule
    \end{tabularx}
\end{table}

\subsubsection{对比模型}
为评估不同方法在当前基准上的长链条平面几何推理能力，本工作选取了三类具有代表性的基线模型：几何专用模型（Geometry Models）、常规非思考型大语言模型（Standard LLMs）和思考型大语言模型（Thinking LLMs）。这一划分主要对应不同的求解范式：几何专用模型通常依赖形式化表示和定制化求解框架，而思考型与常规大语言模型则主要依托通用图文解析与生成能力完成求解。

几何专用模型面向几何任务专门设计，能够较好反映传统几何求解技术的发展水平。例如，Inter-GPS~\cite{lu2021intergps} 将自然语言题干与几何图像解析为形式化语言，并结合定理预测与符号推理完成求解；GeoDRL~\cite{peng2023geodrl} 在逻辑图推理框架中引入强化学习，以优化定理选择策略与搜索路径；Geoformer~\cite{chen2022unigeo} 则将几何计算与证明过程统一建模为序列生成问题。这类方法通常建立在规则系统、符号表示或专用推理框架之上，具有较强的任务针对性。此外，本工作还选用了 G-LLaVA-13B~\cite{gao2023gllava} 与 GeoX~\cite{xia2024geox} 等面向几何任务训练或微调的多模态模型，用于考察几何垂域视觉语言模型在该基准上的实际表现。

通用大模型进一步划分为常规型与思考型两类。类似 GPT-4o~\cite{hurst2024gpt4o}、Claude-4~\cite{anthropic2025claude4} 等常规非思考型大语言模型虽然具备较强的图文理解与生成能力，但并未针对长链推理过程进行显式强化，对其进行测试有助于分析通用模型在复杂几何任务中的原生推理稳定性。思考型模型如 o4-mini~\cite{openai2025o3o4mini} 和 Gemini-2.5-Pro~\cite{comanici2025gemini25}，能够在生成过程中显式展开中间推理过程，因此在多步逻辑推导任务中通常表现更强，可作为观察大语言模型几何推理能力上限的重要参照。

\xsubsection{实验结果与分析}{Baseline Results and Analysis}

\subsubsection{对比实验}
表 \ref{tab:model_comparison_final} 详细展示了不同模型在各个推理步数区间以及全集（All Steps）上的答案准确率（$\mathrm{Acc}_{\text{ans}}$）和过程逻辑正确率（$\mathrm{Acc}_{\text{step}}$）。
% 表格中 1-4 步的区间主要对应了前文所述的短步骤对照组，而 5-12步及以上的区间则重点反映了模型在长链条题目上的表现。
% 1-4steps 509题 5-12steps 1333题 13-26steps 255题
\begin{table}[H]
    \centering
    \caption{不同模型在推理步数区间下的性能对比}
    \label{tab:model_comparison_final}
    \small
    \renewcommand{\arraystretch}{0.8}
    \begin{tabular*}{\textwidth}{@{\extracolsep{\fill}}l cc cc cc cc}
        \toprule
        \multirow{2}{*}{模型} & \multicolumn{2}{c}{1 -- 4 Steps} & \multicolumn{2}{c}{5 -- 12 Steps} & \multicolumn{2}{c}{13 -- 26 Steps} & \multicolumn{2}{c}{All Steps} \\
        \cmidrule(lr){2-3} \cmidrule(lr){4-5} \cmidrule(lr){6-7} \cmidrule(lr){8-9}
        & $\mathrm{Acc}_{\text{ans}}$ & $\mathrm{Acc}_{\text{step}}$ & $\mathrm{Acc}_{\text{ans}}$ & $\mathrm{Acc}_{\text{step}}$ & $\mathrm{Acc}_{\text{ans}}$ & $\mathrm{Acc}_{\text{step}}$ & $\mathrm{Acc}_{\text{ans}}$ & $\mathrm{Acc}_{\text{step}}$ \\
        \midrule
        \multicolumn{9}{@{}l}{\textit{Geometry Models}} \\
        Inter-GPS~\cite{lu2021intergps}             & 44.1 & -    & 35.4 & -    & 7.6 & -    & 34.1 & -    \\
        GeoDRL~\cite{peng2023geodrl}               & 45.8 & -    & 37.4 & -    & 8.4 & -    & 35.9 & -    \\
        Geoformer~\cite{chen2022unigeo}            & 49.0 & -    & 40.3 & -    & 8.2 & -    & 38.5 & -    \\
        G-LLaVA-13B~\cite{gao2023gllava}           & 53.2 & 48.9 & \textbf{44.8} & 42.8 & 9.1 & 37.4 & 42.5 & 43.7 \\
        GeoX~\cite{xia2024geox}                    & \textbf{53.2} & \textbf{50.6} & 45.0 & \textbf{44.3} & \textbf{9.7} & \textbf{38.9} & \textbf{42.7} & \textbf{45.2} \\
        \midrule
        \multicolumn{9}{@{}l}{\textit{Standard LLMs}} \\
        GPT-4o~\cite{hurst2024gpt4o}              & 57.7 & 55.0 & 44.1 & 52.2 & 10.2 & 58.8 & 46.1 & 53.3 \\
        InternVL-3-8B~\cite{zhu2025internvl3}      & 72.8 & 63.7 & 46.3 & 48.3 & 10.4 & 43.2 & 51.4 & 49.8 \\
        Claude-4~\cite{anthropic2025claude4}       & 67.9 & 71.8 & 55.1 & 63.8 & 17.7 & \textbf{81.0} & 56.6 & 68.7 \\
        GPT-4.1~\cite{openai2025gpt41}             & 70.2 & 67.2 & \textbf{56.1} & 65.5 & \textbf{18.9} & 74.7 & 57.8 & \textbf{69.7} \\
        Qwen3-VL-8B~\cite{bai2025qwen3vl}          & \textbf{77.5} & \textbf{89.5} & 55.2 & \textbf{76.6} & 16.5 & 42.9 & \textbf{58.4} & 57.3 \\
        \midrule
        \multicolumn{9}{@{}l}{\textit{Thinking LLMs}} \\
        Gemini-2.0-Thinking~\cite{google2024gemini2}  & 89.7 & 91.2 & 77.6 & 89.7 & 40.2 & 82.5 & 78.7 & 87.4 \\
        Qwen3-VL-8B-Thinking~\cite{bai2025qwen3vl} & 89.5 & 93.1 & 79.5 & 95.2 & 62.2 & 93.0 & 80.7 & 94.5 \\
        Claude-4.5-Thinking~\cite{anthropic2025sonnet45} & 89.3 & 95.4 & 82.3 & 85.9 & 51.2 & 81.7 & 81.7 & 85.7 \\
        o4-mini~\cite{openai2025o3o4mini}          & 94.5 & 96.3 & \textbf{93.6} & 92.5 & \textbf{81.1} & 92.3 & \textbf{92.8} & 95.3 \\
        Gemini-2.5-Pro~\cite{comanici2025gemini25} & \textbf{95.3} & \textbf{96.4} & 91.3 & \textbf{96.7} & 71.3 & \textbf{96.4} & 91.2 & \textbf{96.6} \\
        \bottomrule
    \end{tabular*}
\end{table}

% \xsubsection{结论分析}{Key Findings Analysis}

如表 \ref{tab:model_comparison_final} 所示，实验结果呈现出显著的性能分层格局：思考型大模型展现出压倒性优势，标准大语言模型次之，而几何专用模型表现相对薄弱。这种性能差异不仅体现在最终答案的准确率上，在过程逻辑的严密性上也得到了验证。以全集结果为例，o4-mini 与 Gemini-2.5-Pro 的答案准确率分别为 92.8\% 和 91.2\%，其过程逻辑正确率也高达 95.3\% 与 96.6\%，显著领先于其他类别。相比之下，标准大语言模型虽然在短步骤题目中表现出一定的几何推理潜力，但在复杂场景下仍与思考型模型存在较大差距，而几何专用模型的全集答案准确率普遍受限于 40\% 左右。上述结果表明，几何推理能力的突破不仅依赖于领域知识的引入，更取决于模型在长程推理中维持逻辑一致性与中间状态的能力。

从答案准确率与过程逻辑正确率的关系可以发现，单一的答案评估指标不足以全面刻画模型的真实推理水平。实验观察到两类典型现象：一是部分模型虽然能够获得正确答案，但过程逻辑正确率较低，即通过不稳定的中间推理“偶然”得出正确结论；二是部分模型在复杂题项中虽能保持较高的步骤正确率，但由于长链条推理中的误差累积，最终答案的正确率并未同步提升。例如，Qwen3-VL-8B 在 1--4 步区间下的答案准确率为 77.5\%，但其全集过程逻辑正确率仅为 57.3\%；而 Gemini-2.5-Pro 虽然在全集答案准确率上略低于 o4-mini（91.2\% 和 92.8\%），但其过程逻辑正确率96.6\%显著高于后者。因此，本章引入双重评估维度，以区分“结果偶然正确”与“逻辑相对可靠”这两类不同的模型行为，并为几何推理能力分析提供一定的解释依据。

值得注意的是，几何专用模型在本基准测试中并未展现出预期中的领域优势。尽管部分模型在实验中采用了人工标注的形式化语言输入，规避了视觉理解与文本解析带来的干扰，其整体表现仍大幅落后于通用大模型，且在长链条区间内性能衰减极快。
% 例如，GeoDRL 在 13--26 步区间上的答案准确率仅为 8.4\%，其余专用模型亦处于相近水平。
这表明，即便在理想的输入条件下，传统几何求解器在多条件整合、跨步骤信息保持等方面仍存在明显的瓶颈。
% 本基准不仅量化了各模型的最终求解效能，更揭示了传统几何模型在应对复杂真实场景时的能力上限，充分体现了评测任务的区分度与分析价值。
本基准不仅量化了各模型的最终求解效能，也揭示了传统几何模型在复杂真实场景中的能力上限。

从难度分层的趋势看，长链条任务构成了当前几何推理研究的核心挑战。随着推理步数从 1--4 步增至 13--26 步，所有模型的性能均出现系统性下滑，但衰减幅度差异显著：几何专用模型在长链条区间几近失效，标准大语言模型性能大幅缩水，而思考型大模型虽保持领先，也难以完全避免长链推理带来的负面影响。实验结果证明，长链条几何题并非仅是搜索空间的线性增加，更是对约束保持、定理检索、误差抑制能力的综合考验。因此，本章构建的数据基准通过长链条划分与过程级评估，能够有效拉开模型间的能力差距，为后续几何推理模型的针对性优化提供科学可信的分析依据。

\subsubsection{案例分析}
% \textcolor{red}{[在此处补充一段具体的定性案例分析（Case Study），例如：在测试集的某道 18 步证明题中，XX 模型在第 5 步未能满足推导正确性，凭空捏造了一条角平分线，导致后续虽然得出了正确答案，但其 $\mathrm{Acc}_{\text{step}}$ 被评估机制判定为极低。该案例直观地展现了本基准引入 $\mathrm{Acc}_{\text{step}}$ 指标以剥离“侥幸正确”现象的有效性。]}

% \xsubsection{数据增强有效性验证}{Validation of Data Augmentation Effectiveness}



% -------------------------------------------------------------------
\xsection{本章小结}{Chapter Summary}

本章围绕面向长链条平面几何推理的标准化评测需求，构建了多模态数据基准 GeoLStep。具体而言，本章从真实中考试卷中筛选并整理了 2097 道高质量平面几何题目，完成了题干、图像、推理步骤及辅助结构信息的系统标注，并设计了从自然语言与几何图像到结构化谓词表示的形式化转换链路。在此基础上，本章进一步建立了同时关注最终答案与过程逻辑的评测协议，并对多类代表性模型进行了系统基线测试。实验结果表明，GeoLStep 对当前主流模型具有较高挑战性，尤其在较长推理链条场景下，模型在多步推导稳定性与过程一致性方面仍存在明显不足。本章工作为后续长链条几何推理增强方法的设计与验证提供了统一的数据基础和评测依据。